\documentclass{article}
\usepackage{graphicx} % Required for inserting images

\usepackage{amsmath}

%-------------------------------
%           MARGINS
%--------------------------------

\usepackage[top=2cm]{geometry}

%-------------------------------
%            HEBREW
%-------------------------------

\usepackage{hebrewdoc}

%-------------------------------
%           HYPERLINKS
%-------------------------------

\usepackage{xcolor}
\usepackage{hyperref}

\hypersetup{
    colorlinks=false,
    urlbordercolor=red,
    linkbordercolor=red,
    pdfborder={0 0 1},
    pdfborderstyle={/S/S/W 1}
}

%-------------------------------
%            LISTS
%-------------------------------

\usepackage{enumitem}

%-------------------------------
%            TITLE
%-------------------------------

\title{דף מידע \\ אלגברה ב' - 01040168 \\ חורף תשפ"ז 2026-2027}
\date{}

\begin{document}

\maketitle

\section*{סגל הקורס}

\begin{description}
\item[מרצה:]
אַלַן לֶאוּ
-
\textenglish{\href{mailto:alanlew@technion.ac.il}{alanlew@technion.ac.il}}
\\
\emph{הרצאות:}
יום א' 10:30-12:20 ויום ד' 10:30-12:20, באולמן 102.
\\
\emph{משרד:}
אמאדו 905.
\\
\emph{שעת קבלה:}
יום א' 12:30 במשרד, או בתיאום מראש.

\item[מתרגל:] 
אַלֶן סוֹרַנִי-רוֹזֵט
-
\textenglish{\href{mailto:alan.srozet@technion.ac.il}{alan.srozet@technion.ac.il}}
\\
\emph{תרגולים:}
יום ג' 12:30-14:30, באולמן 102.
\\
\emph{משרד:}
אמאדו 700.
\\
\emph{שעת קבלה:}
יום ג' 14:30-15:20 במשרד, או בתיאום מראש.

\end{description}

\section*{סילבוס}
\begin{enumerate}
\item מרחבי מכפלה פנימית, המשלים הניצב ותהליך גרם־שמידט. הטלות אורתוגונליות.
\item
אופרטורים על מרחבי מכפלה פנימית:
 האופרטור הצמוד, אופרטורים נורמליים, אופרטורים צמודים לעצמם, אופרטורים אוניטריים ואיזומטריות,
משפט הפירוק הספקטרלי ולכסון אורתוגונלי, אופרטורים חיוביים, פירוק לערכים סינגולריים, ופירוק פולארי.
\item סכומים ישרים, הטלות, תת־מרחבים שמורים וצמצום של אופרטור לינארי.
\item אופרטורים נילפוטנטיים ואינדקס הנילפוטנטיות. צורת ז'ורדן.
\item משפט קיילי המילטון ומשפט הפירוק הפרימרי.
\item תבניות בילינאריות ותבניות ריבועיות. חפיפת מטריצות. משפט סילבסטר.
\item אלגברה מולטילינארית: הדטרמיננטה, המטריצה המצורפת, כלל קרמר, ומכפלות טנזוריות.
\end{enumerate}

\section*{ספרות}
\begin{english}
\begin{itemize}
\item[-] \textenglish{Linear Algebra Done Right - S.~Axler}
\item[-] \textenglish{Algebra (2nd Edition) - M.~Artin}
\item[-] \textenglish{Linear Algebra - K.~Hoffman and R.F.~Kunze}
\end{itemize}
\end{english}

\section*{מבנה הציון}

הציון הסופי בקורס יקבע באופן הבא.
\textbf{
ציון מגן ילקח בחשבון בתנאי שהוא גבוה מציון המבחן, ובתנאי שהציון במבחן הינו עובר.}

\begin{itemize}
\item[-] \textbf{5\% מגן עבור שיעורי בית.}
בערך פעם בשבוע יעלה גיליון תרגילי בית עם ציון בינארי, וממוצע הציונים עבורם כפול 100 יהווה מגן של
5\%.

\item[-] \textbf{5-20\% מגן עבור מיני־בחנים.}
במהלך הסמסטר יתקיימו 4 מיני־בחנים, אשר יתקיימו בזמן ומקום ההרצאה בימי ראשון, בתאריכים ה־22.11, 13.12, 3.1, 24.1. יתכנו שינויים בתאריכים, ובמקרה זה נודיע על כך מראש. כל בוחן יכלול שאלה אחת, אולי עם 2-3 סעיפים, ומשכו יהיה 20 דקות (עם תוספת זמן למי שזכאים לכך). החומר לכל בוחן יהיה החומר הנלמד בקורס מאז הבוחן הקודם.
\textbf{ציון כל בוחן יהיה 5\% מגן עבור הציון הסופי.} \\
לא יהיו מועדי ב' למיני־בחנים. סטודנטים שיפספסו מיני־בחנים מסיבות שמוכרות על ידי הטכניון יהיו זכאים ל%
"שאלות בוחן"
במבחן הסופי, שעבורן ינתן ציון מגן עם אותו משקל כמו הבחנים שפוספסו.

\item[-] \textbf{75-100\% בחינה סופית.}
בחינה סופית שאורכה 3 שעות. מועד א' יהיה ב־8.2 ומועד ב' ב־8.3, ומיקומים ושעות עבור שניהם יפורסמו בהמשך על ידי לימודי הסמכה. הבחינה תכלול את כל החומר שנלמד בקורס במהלך הסמסטר, כולל ההרצאות, התרגולים ותרגילי הבית. ציון המבחן הסופי יהווה 75-100\% מהציון בקורס, בהתאם לציוני המגן.
\textbf{שימו לב כי ציוני המגן על תרגילי הבית והמיני־בחנים ילקחו בחשבון רק במקרה של ציון עבור במבחן הסופי (ציון של 55 ומעלה).}
\end{itemize}

\section*{מידע כללי, מילואים והתאמות}

\begin{itemize}
\item[-]
פניות בנושאי ניהול הקורס ושיעורי בית יש להפנות לאלן המתרגל האחראי.
\item[-]
סטודנטים שיהיו במילואים במשך זמן משמעותי או כזה המצריך התייחסות מיוחדת מתבקשים לידע את אלן המתרגל על כך. סטודנטים מקבוצות מילואים 2+3 (המוגדרות על ידי הטכניון
\href{https://ugportal.technion.ac.il/\%D7\%9E\%D7\%99\%D7\%93\%D7\%A2-\%D7\%9C\%D7\%A1\%D7\%98\%D7\%95\%D7\%93\%D7\%A0\%D7\%98\%D7\%99\%D7\%9D/\%D7\%9E\%D7\%99\%D7\%9C\%D7\%95\%D7\%90\%D7\%99\%D7\%9D-\%D7\%97\%D7\%93\%D7\%A9/}{באתר לימודי הסמכה}
)
יהיו זכאים לשימוש בדף נוסחאות במבחן הסופי. דף זה יהיה חד־צדדי ועל כל סטודנט יהיה להכין את דף הנוסחאות שלו בעצמו.
\item[-] בנוגע לזכאות להתאמות מעבר לתוספת זמן, יש לכתוב בנושא לאלן המתרגל כדי שנדון בהן ונקבע איך ליישם אותן במהלך הקורס או בבחינות.
\item[-] \textbf{יש לעבור על ההנחיות עבור הגשת שיעורי הבית שבהמשך מסמך זה.}
\end{itemize}

\section*{נהלי הגשת שיעורי בית}

\begin{enumerate}
\item מותר להגיש תרגילי בית עד 5 ימי איחור מצטברים בלי אישור מיוחד, כאשר איחור של $n$ ימים בהגשה יחושב בתור איחור של
$\lceil n \rceil$
(ערך עליון של
$n$)
ימים, והאיחור המצטבר יהיה סכום האיחורים של ההגשות. הגשות באיחור מעבר לכך לא יבדקו או יחשבו לציון.
\textbf{
הארכות נוספות ינתנו עקב מילואים או מקרים רפואיים חריגים בלבד, באישור בכתב ובתאם לשיקול הדעת של סגל הקורס.
}
\item כל הגשה לא ריקה תקבל ציון מלא
$1$.
כדי לקבל משוב עבור ההגשות, יש להגיש פתרונות קריאים בהתאם לקריטריונים הבאים. הגשות שלא עונות על כך עשויות לא לקבל משוב.
\item על כל הגשה לכלול קובץ
\textenglish{PDF}
יחיד בלבד.
\item את הגיליונות יש להגיש עם גודל עמוד באחד מהגדלים הבאים בלבד, כאשר התוכן עצמו נפרש לאורך מרבית העמוד:
\begin{english}
\begin{itemize}
\item 210 x 297 mm (Portrait A4)
\item 216 x 279 mm (Portrait Letter / ANSI A)
\end{itemize}
\end{english}
אפשר לבדוק את גודל הדף בעורך
\textenglish{PDF}
כמו
\textenglish{Adobe Acrobat Reader}
או
\textenglish{Okular},
ואפשר ברוב התוכנות לטאבלטים לבחור גודל דף מתאים.

\item אין להעלות קובץ הכולל תמונות של הסריקה שאינן מכסות את כל גודל הדף.

\item על הכתב להיות גדול באופן קריא ולכסות לפחות את כל רוחב השורה בדף.
\end{enumerate}

\end{document}